%!TEX root = ../main.tex

\chapter{Θεωρητικό Υπόβαθρο}

\section{Παραλληλοποίηση με Εργασίες}

Στον χώρο του \textbf{HPC}, η παραλληλοποίηση με εργασίες (task parallelism) είναι μια προσέγγιση, μία από τις πλέον διαδεδομένες, που επιτρέπει την εκτέλεση πολλαπλών εργασιών ταυτόχρονα, όπου κάθε εργασία μπορεί να είναι ανεξάρτητη ή να έχει εξαρτήσεις από άλλες εργασίες. Αυτή η μέθοδος είναι ιδιαίτερα χρήσιμη όταν οι εργασίες έχουν διαφορετικά μεγέθη ή απαιτούν διαφορετικούς πόρους. Αυτό έχει ως αποτέλεσμα την καλύτερη εκμετάλλευση των διαθέσιμων πόρων του συστήματος και τη μείωση του χρόνου εκτέλεσης.

Η μετατροπή ενός αλγορίθμου για εκτέλεση σε ένα περιβάλλον παραλληλοποίησης με εργασίες απαιτεί διάφορα βήματα. Αρχικά, η υλοποίηση και η μελέτη του αλγορίθμου σε σειριακή μορφή σε μια γλώσσα προγραμματισμού θα βοηθήσει στην κατανόηση της λειτουργίας του και στην αναγνώριση των τμημάτων που μπορούν να παραλληλοποιηθούν. Ιδανικά αυτό θα θέλαμε να γίνεται αυτόματα από τον μεταγλωττιστή και το λειτουργικό σύστημα, αλλά στην πράξη απαιτείται η παρέμβαση του προγραμματιστή για να καθορίσει τις εργασίες και τις εξαρτήσεις τους \cite{papadakis2012}.

Τα βήματα που πρέπει να λάβει υπόψη ο προγραμματιστής στη διαδικασία παραλληλοποίησης ενός αλγορίθμου, είναι τα εξής \cite{papadakis2012}:
\begin{enumerate}
    \item \textbf{Ανάλυση του Αλγορίθμου σε Εργασίες (tasks):} Ο αλγόριθμος πρέπει να αναλυθεί για να εντοπιστούν τα τμήματα που μπορούν να εκτελεστούν ανεξάρτητα. Αυτό μπορεί να περιλαμβάνει την ανάλυση των βρόχων, των συναρτήσεων και των δεδομένων που χρησιμοποιούνται.
    \item \textbf{Ανάθεση σε Διεργασίες (processes) ή Νήματα (threads):} Οι εργασίες πρέπει να ανατεθούν σε διαφορετικές διεργασίες (ή νήματα) ή να εκτελεστούν σε διαφορετικούς πόρους, με βάση τη διαθεσιμότητα και τις απαιτήσεις του κάθε ενός.
    \item \textbf{Χρονοδρομολόγηση των Διεργασιών (scheduling):} Η χρονοδρομολόγηση των εργασιών είναι ένα σύνθετο βήμα που αναλύεται στα εξής:
    \begin{enumerate}
        \item \label{task_scheduling} \textbf{Ανάθεση κάθε Διεργασίας σε ένα Πόρο:} Κάθε διεργασία πρέπει να ανατεθεί σε έναν πόρο, όπως ένας επεξεργαστής ή ένας πυρήνας, με βάση τις απαιτήσεις της εργασίας και τη διαθεσιμότητα των πόρων.
        \item \textbf{Απόφαση για την Εκτέλεση των Εργασιών:} Δηλαδή, πότε θα εκτελεστεί κάθε εργασία (task) στον συγκεκριμένο πόρο που του αντιστοιχίστηκε στο βήμα \ref{task_scheduling}.
    \end{enumerate}
\end{enumerate}

Τα παραπάνω βήματα παραλληλοποίησης ενός αλγορίθμου με εργασίες φαίνονται και στο παρακάτω Σχήμα \ref{fig:task_parallelization}.

\begin{figure}[h]
    \centering
    \includegraphics[width=0.8\textwidth]{figures/Parallelization-with-tasks-papadakis2012}
    \caption{Βήματα παραλληλοποίησης ενός αλγορίθμου με εργασίες. (Πηγή: \cite{papadakis2012})}
    \label{fig:task_parallelization}
\end{figure}

\section{Χρονοδρομολόγηση Εργασιών}

Όπως είδαμε στην προηγούμενη ενότητα, η χρονοδρομολόγηση των εργασιών είναι το πιο σύνθετο βήμα από τα υπόλοιπα, το οποίο στη σημερινή εποχή γίνεται όλο και πιο δύσκολο λόγω της αυξανόμενης πολυπλοκότητας των συστημάτων HPC και των εφαρμογών που εκτελούνται σε αυτά.

Η χρονοδρομολόγηση των εργασιών περιλαμβάνει την ανάθεση κάθε διεργασίας σε έναν πόρο, όπως ένας επεξεργαστής ή ένας πυρήνας, με βάση τις απαιτήσεις της εργασίας και τη διαθεσιμότητα των πόρων. Επιπλέον, πρέπει να ληφθεί υπόψη η απόφαση για το πότε θα εκτελεστεί κάθε εργασία (task) στον συγκεκριμένο πόρο που του αντιστοιχίστηκε. Αυτό καθιστά την επιλογή μιας αποτελεσματικής στρατηγικής (ή πολιτικής) χρονοδρομολόγησης κρίσιμη για την απόδοση των εφαρμογών σε περιβάλλοντα HPC.

Παλαιότερα ήταν αρκετό να γίνεται η χρήση απλών πολιτικών χρονοδρομολόγησης, όπως η First-Come-First-Served (FCFS) ή η Round Robin (RR), για να επιτευχθεί ικανοποιητική απόδοση. Ωστόσο, με την αύξηση της πολυπλοκότητας των συστημάτων και των εφαρμογών, αυτές οι απλές πολιτικές δεν είναι πλέον επαρκείς. Σήμερα, απαιτούνται πιο εξελιγμένες στρατηγικές χρονοδρομολόγησης που λαμβάνουν υπόψη παράγοντες όπως οι εξαρτήσεις μεταξύ των εργασιών, η εκτίμηση του χρόνου εκτέλεσης κάθε εργασίας, η δυναμική κατάσταση των πόρων και η τοπικότητα των δεδομένων κάθε εργασίας στη μνήμη, διότι πλέον τα συστήματα είναι πολύπλοκα εισάγοντας έννοιες ετερογένειας (heterogeneity) και πολλαπλών επεξεργαστικών μονάδων, όπως οι αρχιτεκτονικές \textbf{NUMA}.

Οπότε θέλουμε η χρονοδρομολόγηση των εργασιών να γίνεται δυναμικά είτε από το λειτουργικό σύστημα είτε από το περιβάλλον εκτέλεσης (runtime environment) με βάση τις τρέχουσες συνθήκες του συστήματος και των εργασιών. Αυτό μπορεί να επιτευχθεί μέσω της χρήσης τεχνικών όπως \textbf{Work Stealing}, \textbf{HEFT} κ.α. που λαμβάνουν υπόψη τους διάφορους παράγοντες που αναφέρθηκαν προηγουμένως. Τέτοιες τεχνικές υλοποιούνται όπως θα δούμε σε παρακάτω κεφάλαια και από τη \textbf{torcpy} και τη \textbf{StarPU} \cite{hadjidoukas2020torcpy, augonnet2009starpu} στο επίπεδο του περιβάλλοντος εκτέλεσης.

\section{Αρχιτεκτονική NUMA}

Όπως αναφέρθηκε και στην προηγούμενη ενότητα, η πολυπλοκότητα των συστημάτων \textbf{HPC} έχει αυξηθεί λόγω της διαρκούς εξέλιξής τους και της ανάγκης για επίλυση πιο σύνθετων επιστημονικών και βιομηχανικών προβλημάτων. Ιστορικά, ωστόσο, το μεγαλύτερο εμπόδιο στην επίτευξη της μέγιστης δυνατής υπολογιστικής απόδοσης δεν ήταν η ταχύτητα των ίδιων των επεξεργαστών, αλλά ο περιορισμός που είναι ευρέως γνωστός ως το «εμπόδιο του Von Neumann» (Von Neumann bottleneck). Το φαινόμενο αυτό περιγράφει την εγγενή ανισορροπία μεταξύ της ταχύτητας επεξεργασίας των Κεντρικών Μονάδων Επεξεργασίας \textbf{(Central Processing Unit - CPU)} και της ταχύτητας μεταφοράς των δεδομένων από και προς την κύρια μνήμη του συστήματος, γεγονός που συχνά αφήνει τους επεξεργαστές σε κατάσταση αναμονής, στερημένους από τα απαραίτητα δεδομένα για την εκτέλεση των υπολογισμών \cite{numa}. 

Αρχικά, για να αντιμετωπιστεί αυτό το πρόβλημα, οι σχεδιαστές συστημάτων υιοθέτησαν μια αρχιτεκτονική που ονομάζεται Ομοιόμορφη Πρόσβαση Μνήμης \textbf{(Uniform Memory Access - UMA)}, η οποία είναι γνωστή και ως Συμμετρική Πολυεπεξεργασία \textbf{(Symmetric Multi-Processing - SMP)}, όπου πολλοί επεξεργαστές μοιράζονται την ίδια φυσική μνήμη, όπως φαίνεται και στο Σχήμα \ref{fig:uma_smp}. Το βασικό πλεονέκτημα αυτής της προσέγγισης ήταν η απλότητά της από πλευράς προγραμματισμού. Όλοι οι επεξεργαστές λάμβαναν απολύτως ισότιμη και ομοιόμορφη πρόσβαση (ίσο χρόνο απόκρισης) σε οποιαδήποτε θέση της μνήμης. Ωστόσο, καθώς ο αριθμός των επεξεργαστών σε ένα σύστημα αυξανόταν, η συνολική απόδοση του συστήματος άρχισε να υποφέρει λόγω της αυξημένης συμφόρησης στη μνήμη \cite{denneman2016numa}.

\begin{figure}[H]
    \centering
    \includegraphics[width=0.8\textwidth]{figures/uma-smp}
    \caption{Αρχιτεκτονική Ομοιόμορφης Πρόσβασης Μνήμης (UMA) ή Συμμετρικής Πολυεπεξεργασίας (SMP). (Πηγή: \cite{denneman2016numa})}
    \label{fig:uma_smp}
\end{figure}

Για την επίλυση του παραπάνω προβλήματος αναπτύχθηκε μια νέα αρχιτεκτονική που ονομάζεται Μη Ομοιόμορφη Πρόσβαση Μνήμης \textbf{(Non-Uniform Memory Access - NUMA)}, όπου κάθε επεξεργαστής έχει τη δική του τοπική μνήμη, αλλά μπορεί επίσης να έχει πρόσβαση στη μνήμη άλλων επεξεργαστών. Παράδειγμα αυτής της αρχιτεκτονικής φαίνεται στο Σχήμα \ref{fig:numa}. Σε αυτή την αρχιτεκτονική, η πρόσβαση στη τοπική μνήμη ενός επεξεργαστή είναι ταχύτερη από την πρόσβαση στη μνήμη άλλων επεξεργαστών, γεγονός που εισάγει μια μη ομοιόμορφη συμπεριφορά στην πρόσβαση στη μνήμη. Η αρχιτεκτονική \textbf{NUMA} επιτρέπει την καλύτερη κλιμάκωση των συστημάτων με πολλούς επεξεργαστές, αλλά απαιτεί επίσης πιο περίπλοκες στρατηγικές χρονοδρομολόγησης και διαχείρισης μνήμης για να μεγιστοποιηθεί η απόδοση των εφαρμογών που εκτελούνται σε αυτά τα συστήματα \cite{denneman2016numa, supermicro_numa}.

\begin{figure}[h]
    \centering
    \includegraphics[width=0.8\textwidth]{figures/numa}
    \caption{Αρχιτεκτονική Μη Ομοιόμορφης Πρόσβασης Μνήμης (NUMA). (Πηγή: \cite{denneman2016numa})}
    \label{fig:numa}
\end{figure}

\section{Προκλήσεις της Αρχιτεκτονικής NUMA στην παραλληλοποίηση με Εργασίες}

Η φαινομενική απλότητα που προσφέρει η λογική αφαίρεση του μοντέλου των εργασιών έρχεται σε άμεση σύγκρουση με την εξαιρετικά πολύπλοκη φυσική πραγματικότητα του υλικού \textbf{NUMA}. Η εγγενής απόκρυψη των αρχιτεκτονικών λεπτομερειών από τον προγραμματιστή σημαίνει ότι το σύστημα χρόνου εκτέλεσης (runtime system) καλείται να λάβει κρίσιμες αποφάσεις σχετικά με την τοποθέτηση δεδομένων και τη δρομολόγηση των εργασιών, χωρίς πάντα να διαθέτει την πλήρη εικόνα της υποκείμενης τοπολογίας. Όταν τα αυτόματα συστήματα δρομολόγησης κατανέμουν εργασίες χωρίς επίγνωση της φυσικής τοποθεσίας των δεδομένων \textbf{(NUMA-oblivious scheduling)}, οι επιπτώσεις στην απόδοση είναι καταστροφικές. Αυτή η αναντιστοιχία οδηγεί σε εκθετική αύξηση της κίνησης στους διαύλους διασύνδεσης (interconnect traffic), ακραία υποβάθμιση του εύρους ζώνης και αδυναμία γραμμικής κλιμάκωσης, διαμορφώνοντας ένα τοπίο στο οποίο η συγκριτική ανάλυση και η πειραματική αξιολόγηση των σύγχρονων συστημάτων παράλληλου προγραμματισμού καθίσταται θεμελιώδους σημασίας για την ακαδημαϊκή έρευνα \cite{Deters_2018}.

Προβλήματα όπως τα παραπάνω καλούνται να λύσουν διάφορες τεχνικές που έχουν αναπτυχθεί για την βελτιστοποίηση της απόδοσης σε συστήματα \textbf{NUMA}. Μία τέτοια τεχνική είναι η \textbf{NUMA-aware Work Stealing (NUMA-WS)}, η οποία τροποποίησε τον πυρήνα του Intel Cilk Plus εφαρμόζοντας ολιστικά την αρχή του \textbf{Work-First} \footnote{Η αρχή του \textbf{Work-First} υποστηρίζει ότι νέες εργασίες πρέπει να εκτελούνται το συντομότερο δυνατόν.}. Αντί να βασίζεται σε τεχνικές τυχαίας κλοπής που οδηγούν σε μεταναστεύσεις μεγάλου κόστους, το \textbf{NUMA-WS} προτείνει τη μηχανική ώθησης (work pushing). Το σύστημα αναθέτει στον εργάτη την προσπάθεια τοποθέτησης ενός υπολογιστικού πλαισίου (full frame) ακριβώς δίπλα στα δεδομένα του, εκβιάζοντας τη δημιουργία τοπικότητας. Το \textbf{NUMA-WS} για να αποφύγει να επιβαρύνει πολύ το τρέχον νήμα-εργάτη εφαρμόζει ώθηση δεδομένων μόνο σε τρία συγκεκριμένα σενάρια: κατά την εκτέλεση πράξεων συγχρονισμού, κατά την επιστροφή θυγατρικού νήματος σε αδρανές γονικό, και μετά από επιτυχή κλοπή εργασίας. Με αυτόν τον τρόπο, σύμφωνα με τη φιλοσοφία \textbf{Work-First}, το κόστος της βελτίωσης της τοπικότητας \textbf{NUMA} δεν επιβαρύνει ποτέ το ενεργό νήμα-εργάτη, αλλά αποσβένεται αποκλειστικά στον αδρανή νήμα-κλέφτη. Το «κατώφλι ώθησης» εγγυάται ότι οι συνολικές λειτουργίες παραμένουν σταθερές και ελεγχόμενες, αποτρέποντας τεράστιο φόρτο εργασίας. Στην πράξη, τα αποτελέσματα είναι εντυπωσιακά: ο πληθωρισμός έργου \textbf{Work Inflation} \footnote{Το \textbf{Work Inflation} αναφέρεται στο φαινόμενο όπου ο συνολικός χρόνος εκτέλεσης που παράγουν όλα τα νήματα μαζί είναι μεγαλύτερος από την εργασία που θα εκτελούσε ένας μόνο επεξεργαστής σειριακά.} μπορεί να μειωθεί σημαντικά στο \textbf{NUMA-WS}, με σημαντική βελτίωση της κλιμάκωσης σε συστήματα με 32 και πλέον πυρήνες \cite{Deters_2018}.

Παρόμοιες τεχνικές έχουν αναπτυχθεί και για άλλες πλατφόρμες, όπως η \textbf{StarPU} \cite{augonnet2009starpu} και η \textbf{torcpy} \cite{hadjidoukas2020torcpy} που θα μελετηθούν στην παρούσα εργασία, οι οποίες υλοποιούν διάφορες στρατηγικές χρονοδρομολόγησης που λαμβάνουν υπόψη την τοπολογία \textbf{NUMA} για να βελτιώσουν την απόδοση των εφαρμογών που εκτελούνται σε τέτοια συστήματα.

\section{Παραλληλοποίηση στην Python}

Η \textbf{Python} είναι μια από τις πιο δημοφιλείς υψηλού επιπέδου γλώσσες προγραμματισμού σήμερα, γνωστή για την απλότητα, την ευκολία χρήσης της και το πλούσιο οικοσύστημα βιβλιοθηκών της. 

Στην καρδιά της \textbf{Python (CPython)} βρίσκεται το \textbf{Global Interpreter Lock (GIL)}, το οποίο είναι ένας μηχανισμός που επιτρέπει μόνο σε ένα νήμα λειτουργικού συστήματος να εκτελεί \textbf{Python bytecode} τη φορά, ακόμα και σε συστήματα με πολλούς επεξεργαστές. Αυτό σημαίνει ότι η παραλληλοποίηση με νήματα (threading) στην \textbf{Python} δεν μπορεί να εκμεταλλευτεί πλήρως τις δυνατότητες των πολυπύρηνων συστημάτων, περιορίζοντας την απόδοση των εφαρμογών που απαιτούν έντονη επεξεργασία \cite{rapidfork2023gil, gross2023pep703, ruys2024deep}. Πιο συγκεκριμένα:

\begin{itemize}
    \item \textbf{CPU-bound Εφαρμογές:} Για υπολογιστικά εντατικές εργασίες, το GIL λειτουργεί ως απόλυτο σημείο συμφόρησης, ακυρώνοντας εντελώς τα πλεονεκτήματα των πολυπύρηνων επεξεργαστών NUMA.
    \item \textbf{I/O-bound Εφαρμογές:} Για εργασίες που εξαρτώνται από αιτήματα δικτύου, λειτουργίες αρχείων κ.λ.π., το GIL είναι εξαιρετικά αποδοτικό (τα νήματα απελευθερώνουν το GIL όταν περιμένουν).
\end{itemize}

Ωστόσο, υπάρχουν διάφορες βιβλιοθήκες και εργαλεία που έχουν αναπτυχθεί για να παρακάμψουν τους περιορισμούς του GIL και να επιτρέψουν την παραλληλοποίηση στην \textbf{Python} \cite{adewole2026concurrency}. Μερικά από αυτά περιλαμβάνουν:

\begin{itemize}
    \item \textbf{Multiprocessing Module:} Η βιβλιοθήκη \textbf{multiprocessing} επιτρέπει τη δημιουργία νέων διεργασιών, οι οποίες μπορούν να εκτελούνται παράλληλα και να εκμεταλλεύονται πολλούς επεξεργαστές, παρακάμπτοντας έτσι το GIL. Ωστόσο, η χρήση διεργασιών μπορεί να είναι πιο βαριά σε πόρους και να απαιτεί περισσότερη μνήμη σε σύγκριση με τα νήματα.
    \item \textbf{Threading Module:} Η βιβλιοθήκη \textbf{threading} επιτρέπει τη δημιουργία νημάτων, αλλά λόγω του GIL, δεν είναι αποτελεσματική για CPU-bound εργασίες.
    \item \textbf{Asyncio Module:} Η βιβλιοθήκη \textbf{asyncio} παρέχει υποστήριξη για ασύγχρονο προγραμματισμό, επιτρέποντας την εκτέλεση πολλαπλών εργασιών (tasks) ταυτόχρονα χωρίς να μπλοκάρει το κύριο νήμα, αλλά δεν είναι κατάλληλη για CPU-bound εργασίες.
    \item \textbf{Concurrent.futures Module:} Η βιβλιοθήκη \textbf{concurrent.futures} παρέχει μια υψηλού επιπέδου διεπαφή για την εκτέλεση εργασιών (tasks) σε νήματα ή διεργασίες, επιτρέποντας την εύκολη παραλληλοποίηση, αλλά εξακολουθεί να υπόκειται στους περιορισμούς του GIL. 
\end{itemize}